%% ============================================================
\section{Discussion}
\label{sec:discussion}
%% ============================================================

\subsection{Lessons Learned}

The development of \echoray{} surfaced several engineering insights that have
broader applicability to the design of collaborative web applications.

\textbf{Simplicity scales surprisingly far.}
The broadcast-relay synchronisation model—broadcasting the full serialised
file tree on every save—was chosen for simplicity over correctness. In
practice, project file trees in the target use case (small-to-medium
prototypes) are rarely larger than a few hundred kilobytes, making full
serialisation over WebSocket negligible in bandwidth terms. The lesson is that
operational transformation or CRDTs are necessary only when \emph{concurrent}
edits from multiple users must be merged without user intervention; for
single-driver workflows (one developer editing while others observe), the
simpler model suffices.

\textbf{Type safety pays dividends at scale.}
The uniform use of TypeScript across client and server, combined with Prisma's
generated types and Zod's runtime validation, eliminated an entire category
of integration bugs during development. Parameters that would have failed
silently (incorrect types, missing required fields) were surfaced at compile
time or at the API boundary, not in production logs.

\textbf{Redis token blacklisting is both simple and effective.}
The token blacklist pattern implemented in \echoray{} provides effective
session revocation with minimal code complexity. The key insight is that
Redis's \texttt{SETEX} command naturally aligns with JWT TTLs: setting the
blacklist entry TTL to the token's remaining validity ensures the cache never
accumulates stale entries.

\textbf{JSON blob storage is a pragmatic trade-off.}
Storing the file tree as a JSON blob in PostgreSQL significantly simplifies the
schema and avoids complex file-entity normalisation. However, it prevents
partial updates (every save overwrites the entire tree) and makes it impossible
to query file contents via SQL. For larger projects, a normalised file model
with content-addressable storage (similar to Git's object model) would be more
appropriate.

\subsection{Comparison with Related Systems}

\textbf{vs GitHub Codespaces/VS Live Share.}
These tools provide superior editor fidelity (true LSP integration, full
extension ecosystem) but require either a local IDE installation (Live Share)
or a paid cloud subscription (Codespaces). \echoray{} provides a fully
browser-native experience with no installation requirement and no per-seat
cost, at the expense of editor feature completeness.

\textbf{vs Replit.}
Replit offers AI assistance and real-time collaboration but uses a proprietary
multiplayer engine whose implementation is not publicly disclosed. \echoray{}
provides full source visibility, enabling researchers and practitioners to
study, extend, and adapt the synchronisation and AI-integration mechanisms.

\textbf{vs CRDT-based editors (Yjs/Automerge).}
Systems built on Yjs or Automerge provide mathematically guaranteed
conflict-free merges, enabling truly offline-first collaboration. \echoray{}
does not currently support offline editing; this is an acknowledged
architectural limitation rather than a deliberate design choice, and
integration of a Yjs provider is a natural extension path.

\subsection{Implications for Practice}

\textbf{AI broadcast as a team synchronisation mechanism.}
The \texttt{@ai} broadcast pattern has implications beyond code generation.
When a team member invokes an AI assistant in \echoray{}, all collaborators
simultaneously see (a) the user's natural-language intent and (b) the AI's
structured response. This shared visibility may reduce the \emph{context
asymmetry} identified by Vaithilingam \etal~\cite{vaithilingam2022expectation}
as a barrier to effective AI-assisted team programming.

\textbf{Structured AI output reduces integration friction.}
By constraining Gemini to output a \texttt{fileTree} object compatible with
the WebContainer API, \echoray{} enables one-click execution of AI-generated
code without any manual copy-paste. This tight integration between AI output
schema and execution substrate is a design pattern worth wider adoption.

\textbf{Redis for token blacklisting should be the default, not the exception.}
Many production systems still rely on short-lived JWTs as the sole revocation
mechanism, accepting a window of up to 15--60 minutes during which a
compromised token remains valid. \echoray{}'s Redis blacklist demonstrates
that immediate revocation is straightforward to implement; the marginal
operational cost (one Redis instance) is justified for any application handling
personal data.

\subsection{Limitations and Open Problems}

\begin{enumerate}[leftmargin=*, label=\textbf{L\arabic*.}]
  \item \textbf{No conflict resolution for concurrent edits.}
    Two users editing the same file simultaneously will experience
    last-writer-wins overwrite behaviour. Integrating a CRDT library (Yjs or
    Automerge) would resolve this.

  \item \textbf{No automated tests.}
    The codebase currently lacks unit tests, integration tests, or end-to-end
    tests. This increases the risk of regression when extending the platform
    and limits confidence in the correctness of the security mechanisms.

  \item \textbf{No rate limiting on the AI endpoint.}
    Unlimited AI queries could exhaust the Gemini API quota and incur
    significant cost. Per-user, per-project, and global rate limits should
    be implemented before public deployment.

  \item \textbf{No streaming AI responses.}
    The current implementation waits for the full Gemini response before
    broadcasting. Streaming (Server-Sent Events or WebSocket streaming) would
    reduce perceived latency for long-form code generation.

  \item \textbf{Single-server Socket.io limitation.}
    As discussed in \cref{sec:evaluation}, scaling beyond one server instance
    requires the Redis adapter. This dependency should be documented and
    included in the default deployment configuration.

  \item \textbf{File-tree size unbounded.}
    There is no server-side limit on the size of the \texttt{fileTree} JSON
    blob. A malicious or careless user could store an arbitrarily large payload,
    degrading database performance. A size cap (e.g., 10\,MB) should be
    enforced at the Zod validation layer.
\end{enumerate}
